% !TEX root = document.tex

\chapter{\label{chap:Preface}Preface}
\addcontentsline{toc}{subsection}{Preface}
When I sat behind my laptop writing Minecraft \textit{mods}\footnote{\textit{modifications}, that change or add to the popular sandbox video game, https://www.minecraft.net/} while watching Java tutorials on YouTube ten years ago, I never would have thought I would end up here, pursuing a degree in Computer Science. I enjoyed most, if not all, subjects taught in school, and on the one hand, I might say I could have ended up pursuing almost any other degree. Deciding what I wanted to do after high school was certainly not easy. On the other hand, looking back, the choice for Computer Science was perhaps inevitable. I remember how programming always had something unique: I enjoyed the problem solving aspect, but also the general idea of making such a mysterious machine do what \textit{I} wanted it to do.

Of course, I quickly realized that the field of Computer Science went a lot deeper than just programming when I arrived in Delft. But, in Delft too, I truly enjoyed almost every course. And at the same time, here too, some courses jumped out: courses focused on low-level programming and computer engineering (e.g. Computer Organization), as well as the more high-level, abstract courses such as Algorithm Design, or Concepts of Programming Languages.

When I finished my bachelor's, I had some doubts about whether I also wanted to pursue a master's. The additional years of education, life experience, and some changes in the field, as well as the world in general, meant I suddenly saw a lot of different options. At the same time, I had really enjoyed the last three years. And because I reject that the sunk cost fallacy is a fallacy, I managed to convince myself that one is not really supposed to do a bachelor's without then doing the master's.

And so, I enrolled. When I did, I decided almost immediately that I would be doing a thesis on a topic related to programming languages, because I had grown to enjoy those courses the most. Jeff's listing for this project was the second one I clicked, and it spoke to me immediately.

I think, looking back objectively, my motivation during the master's was definitely worse than it was during the bachelor's, but I certainly cannot blame this project for that. It was very interesting, as were my weekly meetings with Jeff. We rarely finished on time, but I usually had little else to do for the rest of the day. Jeff, of course, did... sorry! And, of course, thanks for all the help! I would also like to thank Jesper for his feedback, and Neil for joining the committee. Also, I would like to thank them all for being willing to work with me throughout the summer.

Of course, I must also thank my parents and stepparents, for all their support throughout the years. They made it possible for me to focus on my studies without having to worry too much about matters such as money, having a place to live, etc. These things feel normal to me, but I know, of course, that they aren't normal for everyone. As such, I am very grateful for them.

And so, we arrive at the end of the master's. Perhaps fittingly, I once again am faced with doubts about what I want to do now. I have decided that I will not be pursuing a career in academics or really even the field of Computer Science in general. It does feel a little strange, closing this chapter of my life, and `never looking back', in a sense. Of course, I know I will look back: I think my education at the TU has been and will always be incredibly valuable, and I can truthfully say I do not regret any of my choices. Besides the effects of just getting older, I think my five years at the TU have changed me as a person, and perhaps in that sense have caused me to now pursue something different. Well, such is life.

\vspace{1cm}
\begin{flushright}
\theauthor{}\\
Baarn, the Netherlands\\
August 19, 2026\\
\end{flushright}
